\chapter{Topological Manifolds and Group Actions}\label{topological-manifolds}
\chapterstatus{complete}

\section{Introduction}\label{topological-manifolds:section-introduction}

A topological manifold is a space that looks locally like $\RR^n$. The definition has three axioms: locally Euclidean, Hausdorff and second countable. This chapter explains why each is there, by
counterexamples showing that none follows from the other two (Section~\ref{topological-manifolds:section-manifolds}), and develops the two main ways of producing manifolds from old ones: quotients by
group actions (Sections~\ref{topological-manifolds:section-groups} to~\ref{topological-manifolds:section-quotients}) and gluing (Section~\ref{topological-manifolds:section-gluing}). Smooth structures come in
Chapter~\ref{smooth-manifolds}, and homological invariants such as the dimension and the orientation in Chapter~\ref{poincare-duality}. References: \cite{munkres-topology}, \cite{hatcher-at},
\cite{lee-smooth-manifolds}.

\noindent\textbf{Prerequisites.} Chapter~\ref{general-topology} (General Topology), Chapter~\ref{homotopy} (Homotopy and Covering Spaces) and Chapter~\ref{groups} (Groups).

\noindent\textbf{Conventions.} Groups are written multiplicatively, with identity element $e$, unless stated otherwise; $\RR^n$ and $\ZZ^n$ are groups under addition. A \emph{neighbourhood} of a point is an open set
containing it.

\section{Topological groups}\label{topological-manifolds:section-groups}

\begin{definition}\label{topological-manifolds:definition-topological-group}
A \emph{topological group} is a group $G$ with a topology such that the multiplication $G \times G \to G$, $(g, h) \mapsto gh$, and the inversion $G \to G$, $g \mapsto g^{-1}$, are continuous. A \emph{homomorphism} of
topological groups is a continuous group homomorphism.
\end{definition}

\begin{example}\label{topological-manifolds:example-topological-groups}
\begin{enumerate}
\item $(\RR^n, +)$, $(\CC, +)$, the multiplicative groups $\RR^\times$, $\RR_{>0}$ and $\CC^\times$, and the circle $S^1 = \{z \in \CC : |z| = 1\}$, with the Euclidean topologies: the operations are polynomial or rational.
\item $\mathrm{GL}_n(\RR)$ and $\mathrm{GL}_n(\CC)$, open subsets of the spaces of matrices $\RR^{n^2}$ and $\CC^{n^2}$ (as the complement of the zero set of $\det$); multiplication is polynomial in the entries and inversion is rational by
Cramer's rule. Their subgroups $\mathrm{SL}_n$, $\mathrm{O}(n) = \{A : A^TA = I\}$, $\mathrm{SO}(n)$, $\mathrm{U}(n) = \{A : \bar A^TA = I\}$ and $\mathrm{SU}(n)$, with the subspace topologies. $\mathrm{O}(n)$ and $\mathrm{U}(n)$ are
compact, being closed and bounded (their columns are unit vectors).
\item Any group with the discrete topology; products of topological groups.
\item $(\QQ, +)$ with the subspace topology, a topological group that is neither discrete nor locally compact.
\end{enumerate}
\end{example}

\begin{proposition}\label{topological-manifolds:proposition-topological-groups}
Let $G$ be a topological group.
\begin{enumerate}
\item For $g \in G$, left translation $h \mapsto gh$, right translation and inversion are homeomorphisms of $G$.
\item An open subgroup $H$ is also closed.
\item Let $U$ be a neighbourhood of $e$ with $U = U^{-1}$. The subgroup generated by $U$ is open. If $G$ is connected, it is $G$.
\item The connected component $G^0$ of $e$ is a closed normal subgroup.
\item $G$ is Hausdorff if and only if $\{e\}$ is closed.
\end{enumerate}
\end{proposition}

\begin{proof}
(1) Left translation by $g$ is continuous, as the composite of $h \mapsto (g, h)$ and multiplication, with continuous inverse left translation by $g^{-1}$; similarly on the right, and inversion is its own inverse. (2) The
complement of $H$ is the union of the cosets $gH$, $g \notin H$, which are open by (1). (3) The subgroup $H$ generated by $U$ is the union of the sets $u_1 \cdots u_kU$ with $u_i \in U$, which are open by (1). By (2) it is
also closed, and nonempty, so $H = G$ if $G$ is connected. (4) $G^0G^0$ is the image of the connected set $G^0 \times G^0$ under multiplication, hence connected, and it contains $e$; so $G^0G^0 \subseteq G^0$.
Similarly $(G^0)^{-1} \subseteq G^0$ and $gG^0g^{-1} \subseteq G^0$ for all $g$. Components are closed. (5) If $G$ is Hausdorff, points are closed. Conversely, if $\{e\}$ is closed, so is the diagonal
$\{(g, h) : gh^{-1} = e\}$, the preimage of $\{e\}$ under the continuous map $(g, h) \mapsto gh^{-1}$. If $g \neq h$, then $(g, h)$ lies in the open complement of the diagonal, so some product neighbourhood $A \times B$ of
$(g, h)$ misses the diagonal, that is, $A \cap B = \emptyset$.
\end{proof}

\begin{counterexample}\label{topological-manifolds:counterexample-non-hausdorff-group}
Any group with the indiscrete topology, whose only open sets are $\emptyset$ and $G$, is a topological group, and it is not Hausdorff if $|G| > 1$. So (5) is a genuine condition. More generally, for any normal subgroup
$N$ the quotient $G/N$ with the quotient topology is a topological group, and it is Hausdorff if and only if $N$ is closed (Exercise~\ref{topological-manifolds:exercise-quotient-group}).
\end{counterexample}

\section{Group actions and orbit spaces}\label{topological-manifolds:section-actions}

\begin{definition}\label{topological-manifolds:definition-action}
A (left) \emph{action} of a topological group $G$ on a space $X$ is a continuous map $G \times X \to X$, $(g, x) \mapsto g \cdot x = gx$, with $ex = x$ and $g(hx) = (gh)x$. The \emph{orbit} of $x$ is $Gx = \{gx : g \in G\}$, the
\emph{stabiliser} is $G_x = \{g : gx = x\}$, and the \emph{orbit space} $X/G$ is the set of orbits with the quotient topology for the projection $\pi \colon X \to X/G$. The action is \emph{free} if all stabilisers are
trivial, and \emph{transitive} if there is only one orbit.
\end{definition}

Each $g \in G$ acts by a homeomorphism of $X$, with inverse the action of $g^{-1}$. The orbits partition $X$, and $gG_xg^{-1} = G_{gx}$.

\begin{proposition}\label{topological-manifolds:proposition-orbit-space}
Let $G$ act on $X$, with $\pi \colon X \to X/G$.
\begin{enumerate}
\item $\pi$ is open: for $U$ open, $\pi^{-1}(\pi(U)) = \bigcup_{g \in G}gU$ is open.
\item $X/G$ is Hausdorff if and only if the \emph{orbit relation} $R = \{(x, gx) : x \in X, g \in G\}$ is closed in $X \times X$.
\item If $X$ is second countable, so is $X/G$; if $X$ is compact, connected or path-connected, so is $X/G$.
\item If $G$ is compact and $X$ Hausdorff, then $X/G$ is Hausdorff, $\pi$ is closed, and for each $x$ the map $G/G_x \to Gx$, $gG_x \mapsto gx$, is a homeomorphism.
\end{enumerate}
\end{proposition}

\begin{proof}
(1) Each $gU$ is open. (2) Suppose $R$ is closed and $\pi(x) \neq \pi(y)$. Then $(x, y) \notin R$, so there are open $U \ni x$ and $V \ni y$ with $(U \times V) \cap R = \emptyset$. The sets $\pi(U)$ and $\pi(V)$ are open by (1) and
disjoint: if $\pi(u) = \pi(v)$ then $(u, v) \in R$. Conversely, if $X/G$ is Hausdorff, its diagonal is closed (a point $(a, b)$ with $a \neq b$ has the neighbourhood $A \times B$, for disjoint neighbourhoods $A \ni a$,
$B \ni b$, which misses the diagonal), and $R$ is its preimage under the continuous map $\pi \times \pi$. (3) The images of a countable basis form a countable basis, since $\pi$ is open, continuous and surjective;
continuous images of compact, connected or path-connected spaces have the same property.

(4) We first show that projections along a compact factor are closed: if $K$ is compact, $Y$ any space and $C \subseteq K \times Y$ closed, then $\mathrm{pr}_Y(C)$ is closed. Indeed, if $y \notin \mathrm{pr}_Y(C)$, each $k \in K$ has
a product neighbourhood $A_k \times B_k$ of $(k, y)$ disjoint from $C$; finitely many $A_{k_1}, \ldots, A_{k_r}$ cover $K$, and $B = \bigcap_jB_{k_j}$ is a neighbourhood of $y$ with $(K \times B) \cap C = \emptyset$, so $B$ misses
$\mathrm{pr}_Y(C)$. Now let $C \subseteq X$ be closed. Then $\pi^{-1}(\pi(C)) = GC$ is the projection to $X$ of $\{(g, x) : g^{-1}x \in C\}$, which is closed in $G \times X$; so $GC$ is closed, and $\pi$ is a closed map. Next, as $X$
is Hausdorff its diagonal $\Delta_X$ is closed (as in (2)), and $R$ is the projection to $X \times X$ of the closed set $\{(g, x, y) : (gx, y) \in \Delta_X\} \subseteq G \times X \times X$; so $R$ is closed and $X/G$ is Hausdorff by (2).
Finally $g \mapsto gx$ induces a continuous bijection from $G/G_x$, which is compact as the image of $G$, to the Hausdorff space $Gx$; it is a homeomorphism by
Proposition~\ref{general-topology:proposition-compact-basic}(4).
\end{proof}

\begin{example}\label{topological-manifolds:example-orbit-spaces}
\begin{enumerate}
\item $\ZZ$ acts on $\RR$ by translations. The map $t \mapsto e^{2\pi it}$ induces a continuous bijection $\RR/\ZZ \to S^1$, which is a homeomorphism because $\RR/\ZZ$ is compact, the image of $[0, 1]$. Likewise
$\RR^n/\ZZ^n \cong (S^1)^n = T^n$, the $n$-torus.
\item $\{\pm1\}$ acts on $S^n$ by $x \mapsto \pm x$, and $S^n/\{\pm1\} = \RR\PP^n$, the real projective space.
\item $\RR^\times$ acts on $\RR^{n+1} \setminus \{0\}$ by scalar multiplication, and the orbit space is again $\RR\PP^n$; $\CC^\times$ acting on $\CC^{n+1} \setminus \{0\}$ gives the complex projective space $\CC\PP^n$, and $S^1$
acting on $S^{2n+1} \subseteq \CC^{n+1}$ gives $\CC\PP^n$ as well (Proposition~\ref{topological-manifolds:proposition-projective-spaces}).
\item $\mathrm{O}(n)$ acts transitively on $S^{n-1}$, with the stabiliser of $e_n$ equal to $\mathrm{O}(n - 1)$ (embedded in the upper left corner). By Proposition~\ref{topological-manifolds:proposition-orbit-space}(4),
$S^{n-1} \cong \mathrm{O}(n)/\mathrm{O}(n - 1)$; similarly $S^{2n-1} \cong \mathrm{U}(n)/\mathrm{U}(n - 1)$.
\end{enumerate}
\end{example}

\begin{lemma}[Kronecker]\label{topological-manifolds:lemma-kronecker}
If $\alpha \in \RR$ is irrational, the set $\{n\alpha - \lfloor n\alpha\rfloor : n \in \ZZ\}$ of fractional parts is dense in $[0, 1]$.
\end{lemma}

\begin{proof}
Let $N \geq 1$. The $N + 1$ fractional parts of $0, \alpha, \ldots, N\alpha$ lie in $[0, 1)$, which is the union of the $N$ intervals $[j/N, (j+1)/N)$, so two of them, of $k\alpha$ and $l\alpha$ with $k < l$, lie in the same
interval. Then $m = l - k$ satisfies $0 < |m\alpha - p| < 1/N$ for some integer $p$; the first inequality holds because $\alpha$ is irrational. Replacing $m$, $p$ by $-m$, $-p$ if necessary, we may assume that
$\beta = m\alpha - p$ satisfies $0 < \beta < 1/N$. The numbers $j\beta$, $j = 0, 1, 2, \ldots$, increase in steps of $\beta < 1/N$, so every interval $[c, c + 1/N) \subseteq [0, 1)$ contains some $j\beta$; and for $0 \leq j\beta < 1$,
$j\beta = (jm)\alpha - jp$ is the fractional part of $(jm)\alpha$. As $N$ is arbitrary, the set is dense.
\end{proof}

\begin{counterexample}\label{topological-manifolds:counterexample-irrational-rotation}
Orbit spaces can be very bad. Let $\alpha$ be irrational and let $\ZZ$ act on $S^1$ by $n \cdot z = e^{2\pi in\alpha}z$. By Lemma~\ref{topological-manifolds:lemma-kronecker} every orbit is dense. If $U \subseteq S^1$ is open, saturated and
nonempty, then for every $y \in S^1$ the orbit of $y$ meets $U$, so $y \in U$; hence $U = S^1$. The orbit space $S^1/\ZZ$ is therefore an uncountable set with the indiscrete topology: it is not Hausdorff, and every map from it to a
Hausdorff space is constant. The action is free.
\end{counterexample}

\section{Covering space actions and quotients}\label{topological-manifolds:section-quotients}

From now on in this section $G$ is a group acting on $X$ by homeomorphisms, that is, a topological group with the discrete topology.

\begin{definition}\label{topological-manifolds:definition-covering-action}
An action of a discrete group $G$ on $X$ is
\begin{enumerate}
\item a \emph{covering space action} if every $x \in X$ has a neighbourhood $U$ with $gU \cap U = \emptyset$ for all $g \neq e$;
\item \emph{proper} if for every compact $K \subseteq X$ the set $\{g \in G : gK \cap K \neq \emptyset\}$ is finite.
\end{enumerate}
\end{definition}

A covering space action is free. An action of a finite group is proper.

\begin{theorem}\label{topological-manifolds:theorem-covering-action}
Let $G$ act on $X$ by a covering space action. Then $\pi \colon X \to X/G$ is a covering map. If $X$ is path-connected, then $\pi$ is a normal covering whose deck transformation group is $G$, and
$\pi_1(X/G)/\pi_*\pi_1(X) \cong G$; in particular $\pi_1(X/G) \cong G$ if $X$ is simply connected.
\end{theorem}

\begin{proof}
Let $x \in X$ and $U$ as in Definition~\ref{topological-manifolds:definition-covering-action}. Then $\pi^{-1}(\pi(U)) = \bigsqcup_{g \in G}gU$, a disjoint union of open sets: if $gU \cap hU \neq \emptyset$ then $h^{-1}gU \cap U \neq \emptyset$, so
$g = h$. The restriction of $\pi$ to $gU$ is continuous, open (Proposition~\ref{topological-manifolds:proposition-orbit-space}(1)) and injective, since $\pi(gu) = \pi(gu')$ means $gu' = hgu$ for some $h$, whence
$g^{-1}hg = e$ and $u = u'$. So $\pi|_{gU}$ is a homeomorphism onto the open set $\pi(U)$, and $\pi(U)$ is evenly covered.

Each $g$ is a deck transformation, since $\pi(gx) = \pi(x)$, and $G$ acts transitively on each fibre $\pi^{-1}(\pi(x)) = Gx$. If $X$ is path-connected, a deck transformation is determined by the image of one point
(Chapter~\ref{homotopy}), so every deck transformation $\tau$ equals the $g$ with $gx = \tau(x)$, and $g \mapsto (x \mapsto gx)$ is an isomorphism $G \to \mathrm{Deck}(\pi)$, injective because the action is free.
The deck group acts transitively on the fibres, so the covering is normal, and Proposition~\ref{homotopy:proposition-deck}(2) gives $\pi_1(X/G)/\pi_*\pi_1(X) \cong \mathrm{Deck}(\pi) \cong G$.
\end{proof}

\begin{theorem}\label{topological-manifolds:theorem-proper-free}
Let a discrete group $G$ act freely and properly on a locally compact Hausdorff space $X$. Then the action is a covering space action and $X/G$ is Hausdorff.
\end{theorem}

\begin{proof}
Let $x \in X$ and let $K$ be a compact neighbourhood of $x$ (containing an open neighbourhood of $x$). The set $F = \{g \neq e : gK \cap K \neq \emptyset\}$ is finite. For $g \in F$, $gx \neq x$ as the action is free, so there are
disjoint open sets $V_g \ni x$ and $W_g \ni gx$. Let $U = \operatorname{int}K \cap \bigcap_{g \in F}(V_g \cap g^{-1}W_g)$, an open neighbourhood of $x$. For $g \neq e$ not in $F$, $gU \cap U \subseteq gK \cap K = \emptyset$; for $g \in F$,
$gU \subseteq W_g$ and $U \subseteq V_g$, which are disjoint. So the action is a covering space action.

For the Hausdorff property we show that the orbit relation $R$ is closed (Proposition~\ref{topological-manifolds:proposition-orbit-space}(2)). Let $(x, y) \notin R$ and choose compact neighbourhoods $K_x$ of $x$ and $K_y$ of $y$. With
$K = K_x \cup K_y$, the set $F = \{g : gK_x \cap K_y \neq \emptyset\}$ is contained in $\{g : gK \cap K \neq \emptyset\}$, hence finite. For $g \in F$, $gx \neq y$, so there are disjoint open $A_g \ni gx$ and $B_g \ni y$. Put
$U = \operatorname{int}K_x \cap \bigcap_{g \in F}g^{-1}A_g$ and $V = \operatorname{int}K_y \cap \bigcap_{g \in F}B_g$. For every $g$, $gU \cap V = \emptyset$: if $g \notin F$ because $gK_x \cap K_y = \emptyset$, and if $g \in F$ because
$gU \subseteq A_g$ and $V \subseteq B_g$. Hence $(U \times V) \cap R = \emptyset$.
\end{proof}

\begin{counterexample}\label{topological-manifolds:counterexample-non-hausdorff-quotient}
A covering space action need not have a Hausdorff quotient. Let $\ZZ$ act on $X = \RR^2 \setminus \{0\}$ by $n \cdot (x, y) = (2^nx, 2^{-n}y)$. It is a covering space action: a point with $x \neq 0$ has the neighbourhood
$\{(x', y') : |x|/\sqrt2 < |x'| < \sqrt2|x|\}$, which is moved off itself by every $n \neq 0$, and similarly if $y \neq 0$. But the points $p = (1, 0)$ and $q = (0, 1)$ lie in different orbits, while
$(1, 2^{-k}) \to p$ and $(-k) \cdot (1, 2^{-k}) = (2^{-k}, 1) \to q$; so $(p, q)$ lies in the closure of the orbit relation, which is not closed, and $X/\ZZ$ is not Hausdorff. The action is free but not proper: the compact set
$K = \{(x, y) : \max(|x|, |y|) = 1\}$ meets its translates by all $n$.
\end{counterexample}

\begin{example}\label{topological-manifolds:example-covering-actions}
\begin{enumerate}
\item $\ZZ^n$ acts freely and properly on $\RR^n$ by translations; $T^n = \RR^n/\ZZ^n$ and $\pi_1(T^n) \cong \ZZ^n$, as $\RR^n$ is simply connected.
\item $\{\pm1\}$ acts freely on $S^n$, so $\pi_1(\RR\PP^n) \cong \ZZ/2$ for $n \geq 2$, since $S^n$ is simply connected for $n \geq 2$ (Chapter~\ref{homotopy}).
\item (Lens spaces) Let $p \geq 2$ and $q$ be coprime integers, $\zeta = e^{2\pi i/p}$. The group $\ZZ/p$ acts freely on $S^3 \subseteq \CC^2$ by $k \cdot (z_1, z_2) = (\zeta^kz_1, \zeta^{qk}z_2)$: if $\zeta^kz_1 = z_1$ and
$\zeta^{qk}z_2 = z_2$ with $(z_1, z_2) \neq 0$, then $\zeta^k = 1$ or $\zeta^{qk} = 1$, and either forces $p \mid k$ as $q$ is prime to $p$. The quotient $L(p, q) = S^3/(\ZZ/p)$ is the \emph{lens space}, with
$\pi_1 \cong \ZZ/p$.
\item (Klein bottle) The transformations $a(x, y) = (x + \frac12, -y)$ and $b(x, y) = (x, y + 1)$ of $\RR^2$ generate a group $G$ acting freely and properly, with $aba^{-1} = b^{-1}$; the quotient $\RR^2/G$ is the Klein bottle, with
$\pi_1 \cong G$, a non-abelian group.
\end{enumerate}
\end{example}

\section{Topological manifolds}\label{topological-manifolds:section-manifolds}

\begin{definition}\label{topological-manifolds:definition-manifold}
Let $n \geq 0$. A space $M$ is \emph{locally Euclidean of dimension $n$} if every point has a neighbourhood homeomorphic to an open subset of $\RR^n$ (equivalently, to $\RR^n$ itself, since every point of an open subset of
$\RR^n$ has a ball neighbourhood, and balls are homeomorphic to $\RR^n$). A \emph{topological manifold of dimension $n$}, or \emph{$n$-manifold}, is a space that is
\begin{enumerate}
\item locally Euclidean of dimension $n$,
\item Hausdorff, and
\item second countable.
\end{enumerate}
A \emph{chart} is a homeomorphism $\varphi \colon U \to \varphi(U)$ from an open $U \subseteq M$ onto an open subset of $\RR^n$.
\end{definition}

That the dimension of a nonempty manifold is determined by the space is a theorem, proved with homology (Chapter~\ref{poincare-duality}). The next three counterexamples show that none of the three axioms follows from the
other two.

\begin{counterexample}[Not Hausdorff]\label{topological-manifolds:counterexample-not-hausdorff}
The \emph{line with two origins} $L$ is the quotient of $\RR \times \{1, 2\}$ by $(t, 1) \sim (t, 2)$ for $t \neq 0$ (Counterexample~\ref{general-topology:counterexample-bad-quotients}). The quotient map is open, the two copies of $\RR$
map homeomorphically onto open subsets covering $L$, so $L$ is locally Euclidean of dimension $1$ and second countable. But the two origins $0_1$, $0_2$ have no disjoint neighbourhoods: every neighbourhood of $0_i$
contains the images of $(-\varepsilon, \varepsilon) \setminus \{0\}$ for some $\varepsilon > 0$. In $L$ the sequence $1/n$ converges to both origins, so limits are not unique.
\end{counterexample}

\begin{counterexample}[Not second countable]\label{topological-manifolds:counterexample-not-second-countable}
Let $I$ be an uncountable set with the discrete topology, and $M = I \times \RR$, a disjoint union of uncountably many lines. It is locally Euclidean of dimension $1$ and Hausdorff. It is not second countable: a basis must
contain, for each $i \in I$, a nonempty open subset of $\{i\} \times \RR$, and these are pairwise distinct. Connected examples also exist: the \emph{long line} and the \emph{Pr\"ufer surface} are connected, Hausdorff and locally
Euclidean of dimensions $1$ and $2$ but not second countable \cite{munkres-topology}. Without second countability, partitions of unity (Theorem~\ref{general-topology:theorem-partitions-of-unity}) and embeddings in
$\RR^N$ can fail.
\end{counterexample}

\begin{lemma}\label{topological-manifolds:lemma-punctured-connected}
Let $n \geq 2$, $W \subseteq \RR^n$ open and connected, and $p \in W$. Then $W \setminus \{p\}$ is connected. For $n = 1$, removing a point from an open interval leaves exactly two components.
\end{lemma}

\begin{proof}
$\RR^n \setminus \{0\}$ is path-connected for $n \geq 2$: two points $a$, $b$ are joined by the segment $[a, b]$ unless $0$ lies on it, and then by the segments $[a, c]$ and $[c, b]$ for any $c$ not on the line through $a$ and $b$.
Hence every punctured ball $B \setminus \{p\}$, homeomorphic to $\RR^n \setminus \{0\}$, is path-connected. Choose a ball $B$ with $p \in B \subseteq W$ and let $A$ be the set of points of $W \setminus \{p\}$ that can be joined to a
point of $B \setminus \{p\}$ by a path in $W \setminus \{p\}$. Every point of $W \setminus \{p\}$ has a ball neighbourhood in $W \setminus \{p\}$, which is path-connected; so $A$ and its complement in $W \setminus \{p\}$ are open. If
the complement were nonempty, then $W = (A \cup B) \cup (W \setminus (A \cup \{p\}))$ would be a decomposition into disjoint nonempty open sets ($A \cup B$ is open, and $B \setminus \{p\} \subseteq A$), contradicting connectedness. So
$W \setminus \{p\} = A$ is path-connected. The case $n = 1$ is clear, as connected subsets of $\RR$ are intervals (Theorem~\ref{general-topology:theorem-intervals-connected}).
\end{proof}

\begin{counterexample}[Not locally Euclidean]\label{topological-manifolds:counterexample-not-locally-euclidean}
The union $X = \{(x, y) \in \RR^2 : xy = 0\}$ of the two coordinate axes is Hausdorff and second countable, being a subspace of $\RR^2$, but not locally Euclidean at the origin. Suppose $\varphi \colon U \to \varphi(U) \subseteq \RR^n$ were a
chart at $0$. Choose $\varepsilon > 0$ with $V = X \cap (-\varepsilon, \varepsilon)^2 \subseteq U$. Then $V$ is a cross with four open arms, open in $X$ and connected, and $V \setminus \{0\}$ has exactly four components, the arms.
So $W = \varphi(V)$ is a connected open subset of $\RR^n$, and $W \setminus \{\varphi(0)\}$, homeomorphic to $V \setminus \{0\}$, has four components. This contradicts Lemma~\ref{topological-manifolds:lemma-punctured-connected} for
$n \geq 1$, and $n = 0$ is impossible because $\{0\}$ is not open in $X$.

Similarly the closed half-line $[0, +\infty)$ is not locally Euclidean at $0$. If $\varphi$ were a chart at $0$ with values in $\RR^n$: $n = 0$ is impossible as before; for $n = 1$, a connected neighbourhood $[0, \varepsilon) \subseteq U$
would map to an open interval, which is disconnected by removing the image of $0$, while $(0, \varepsilon)$ is connected; for $n \geq 2$, a small interval $J \subseteq (0, \varepsilon)$ would map to a connected open subset of $\RR^n$ that
stays connected after removing a point, while $J$ does not. It is a manifold with boundary (Definition~\ref{topological-manifolds:definition-boundary}).
\end{counterexample}

\begin{proposition}\label{topological-manifolds:proposition-manifold-properties}
Let $M$ be an $n$-manifold.
\begin{enumerate}
\item $M$ is locally compact and locally path-connected. Its connected components are open, coincide with its path components, and are countable in number.
\item Every open subset of $M$ is an $n$-manifold.
\item $M$ is paracompact and admits partitions of unity subordinate to any open cover.
\item $M$ has a countable cover by open sets $B_i$ with homeomorphisms $B_i \cong \RR^n$ (\emph{coordinate balls}).
\end{enumerate}
\end{proposition}

\begin{proof}
(1) Each point has a neighbourhood homeomorphic to $\RR^n$, and closed balls and open balls in $\RR^n$ give compact and path-connected neighbourhoods. In a locally path-connected space, path components are open, hence
also closed, and so they are the components. They are disjoint open sets, each containing an element of a countable basis, so there are countably many. (2) The three axioms pass to open subspaces. (3) By (1) and
Hausdorffness $M$ is locally compact Hausdorff and second countable, so Proposition~\ref{general-topology:proposition-exhaustion} and Theorem~\ref{general-topology:theorem-partitions-of-unity} apply. (4) Let
$\mathcal{B}$ be a countable basis and $\mathcal{B}'$ the set of $W \in \mathcal{B}$ contained in the domain of some chart; for each $W \in \mathcal{B}'$ fix a chart $\varphi_W$ whose domain contains $W$. The sets $\varphi_W^{-1}(B)$, for
$W \in \mathcal{B}'$ and $B$ an open ball with rational centre and radius contained in $\varphi_W(W)$, are countably many coordinate balls. They cover $M$: a point $x$ lies in a chart domain, hence in some $W \in \mathcal{B}$ inside it,
so $W \in \mathcal{B}'$, and $\varphi_W(x)$ lies in a rational ball contained in the open set $\varphi_W(W)$.
\end{proof}

\begin{theorem}\label{topological-manifolds:theorem-countable-fundamental-group}
The fundamental group of a connected manifold is countable.
\end{theorem}

\begin{proof}
Let $\{B_i\}$ be a countable cover of $M$ by coordinate balls (Proposition~\ref{topological-manifolds:proposition-manifold-properties}(4)); each $B_i$ is simply connected, being homeomorphic to $\RR^n$. For each pair $i, j$, the open set
$B_i \cap B_j$ has countably many path components (by Proposition~\ref{topological-manifolds:proposition-manifold-properties}(1)); choose a point in each, and let $S$ be the countable set of all chosen points, together with a
base point $p$. For any two points $s, s' \in S$ lying in a common $B_i$, choose a path $\sigma^i_{s,s'}$ in $B_i$ from $s$ to $s'$; these are countably many paths. We claim that every loop $\gamma$ at $p$ is homotopic to a finite
concatenation of such paths; since there are countably many finite concatenations, this proves the theorem.

By the Lebesgue number lemma there is a partition $0 = t_0 < t_1 < \cdots < t_k = 1$ with $\gamma([t_{m-1}, t_m]) \subseteq B_{i_m}$. We may take $p \in B_{i_1} \cap B_{i_k}$. For $0 < m < k$, the point $\gamma(t_m)$ lies in
$B_{i_m} \cap B_{i_{m+1}}$, in a path component containing some $s_m \in S$; choose a path $\tau_m$ in that component from $\gamma(t_m)$ to $s_m$, and let $\tau_0$, $\tau_k$ be constant at $p = s_0 = s_k$. Then $\gamma$ is homotopic to
the concatenation over $m$ of the paths $\bar\tau_{m-1} \cdot \gamma|_{[t_{m-1},t_m]} \cdot \tau_m$, each of which runs in $B_{i_m}$ from $s_{m-1}$ to $s_m$. As $B_{i_m}$ is simply connected, each is homotopic relative to its
endpoints to $\sigma^{i_m}_{s_{m-1},s_m}$.
\end{proof}

\begin{corollary}\label{topological-manifolds:corollary-covering-manifold}
If $M$ is a connected $n$-manifold and $p \colon E \to M$ a covering map with $E$ connected, then $E$ is an $n$-manifold.
\end{corollary}

\begin{proof}
$E$ is locally homeomorphic to $M$, hence locally Euclidean. It is Hausdorff: two points in different fibres are separated by preimages of disjoint neighbourhoods in $M$, and two points in the same fibre lie in different
sheets over an evenly covered neighbourhood. The fibres are in bijection with the cosets of $p_*\pi_1(E)$ in the countable group $\pi_1(M)$ (Chapter~\ref{homotopy}), hence countable. If $\{W_k\}$ is a countable basis of $M$
consisting of evenly covered sets, then the sheets over the $W_k$ form a countable family of open sets, and the sheets over basic open subsets of the $W_k$ give a countable basis of $E$.
\end{proof}

\begin{definition}\label{topological-manifolds:definition-boundary}
Let $\mathbb{H}^n = \{x \in \RR^n : x_n \geq 0\}$. An \emph{$n$-manifold with boundary} is a Hausdorff, second countable space in which every point has a neighbourhood homeomorphic to an open subset of $\mathbb{H}^n$. A point is an
\emph{interior point} if it has a neighbourhood homeomorphic to $\RR^n$, and a \emph{boundary point} otherwise; the boundary is written $\partial M$.
\end{definition}

A point sent by a chart to a point with $x_n > 0$ is interior. That a point sent by some chart to $\{x_n = 0\}$ is not interior, so that $\partial M$ is the set of such points and $\partial M$ is an $(n-1)$-manifold, uses invariance
of domain (Chapter~\ref{poincare-duality}); for $n = 1$ it follows from Lemma~\ref{topological-manifolds:lemma-punctured-connected}.

\begin{example}\label{topological-manifolds:example-boundary}
The closed disc $D^n$ is an $n$-manifold with boundary $S^{n-1}$; $[0, 1]$ has boundary $\{0, 1\}$; the \emph{M\"obius band}, the quotient of $[0, 1] \times [-1, 1]$ by $(0, t) \sim (1, -t)$, has boundary a single circle.
\end{example}

\section{Constructions of manifolds}\label{topological-manifolds:section-constructions}

\begin{proposition}\label{topological-manifolds:proposition-constructions}
\begin{enumerate}
\item Open subsets of $n$-manifolds and countable disjoint unions of $n$-manifolds are $n$-manifolds.
\item If $M$ is an $m$-manifold and $N$ an $n$-manifold, $M \times N$ is an $(m + n)$-manifold.
\item If a discrete group $G$ acts freely and properly on an $n$-manifold $M$, then $M/G$ is an $n$-manifold, and $M \to M/G$ is a covering map.
\item Connected covering spaces of connected $n$-manifolds are $n$-manifolds.
\end{enumerate}
\end{proposition}

\begin{proof}
(1) Clear. (2) Products of charts are charts, $\RR^m \times \RR^n = \RR^{m+n}$, and products of Hausdorff or second countable spaces are Hausdorff or second countable. (3) By
Theorem~\ref{topological-manifolds:theorem-proper-free} and Proposition~\ref{topological-manifolds:proposition-manifold-properties}(1) the quotient is Hausdorff and $\pi$ is a covering map, so $M/G$ is locally homeomorphic to
$M$; it is second countable by Proposition~\ref{topological-manifolds:proposition-orbit-space}(3). (4) is Corollary~\ref{topological-manifolds:corollary-covering-manifold}.
\end{proof}

\begin{example}\label{topological-manifolds:example-manifolds}
\begin{enumerate}
\item $\RR^n$; open subsets such as $\mathrm{GL}_n(\RR) \subseteq \RR^{n^2}$.
\item The sphere $S^n \subseteq \RR^{n+1}$, with the two stereographic charts (Example~\ref{general-topology:example-stereographic}); it is Hausdorff and second countable as a subspace of $\RR^{n+1}$.
\item The graph $\{(x, f(x))\}$ of a continuous map $f \colon U \to \RR^k$, $U \subseteq \RR^n$ open, is an $n$-manifold, homeomorphic to $U$ by the projection.
\item The torus $T^n = \RR^n/\ZZ^n$, $\RR\PP^n = S^n/\{\pm1\}$, the lens spaces and the Klein bottle of Example~\ref{topological-manifolds:example-covering-actions}, by Proposition~\ref{topological-manifolds:proposition-constructions}(3).
\end{enumerate}
\end{example}

\begin{proposition}\label{topological-manifolds:proposition-projective-spaces}
$\CC\PP^n = (\CC^{n+1} \setminus \{0\})/\CC^\times$ is a compact $2n$-manifold, homeomorphic to $S^{2n+1}/S^1$. The sets $U_i = \{[z] : z_i \neq 0\}$ are open and $\varphi_i([z]) = (z_0/z_i, \ldots, \widehat{z_i/z_i}, \ldots, z_n/z_i)$ is a
homeomorphism $U_i \to \CC^n$. The same holds for $\RR\PP^n$, with $\RR$ in place of $\CC$.
\end{proposition}

\begin{proof}
The inclusion $S^{2n+1} \to \CC^{n+1} \setminus \{0\}$ induces a continuous map $S^{2n+1}/S^1 \to \CC\PP^n$, which is bijective since every line meets the sphere in a circle, and has continuous inverse induced by
$z \mapsto z/|z|$. By Proposition~\ref{topological-manifolds:proposition-orbit-space}(3),(4) applied to the compact group $S^1$, $\CC\PP^n$ is compact, Hausdorff and second countable. The set $U_i$ is open because its preimage
$\{z_i \neq 0\}$ is. The map $\varphi_i$ is well defined and continuous (it is induced by a continuous map constant on orbits), with continuous inverse $w \mapsto [w_0 : \cdots : 1 : \cdots : w_n]$, the $1$ in position $i$.
\end{proof}

\begin{example}\label{topological-manifolds:example-cp1-sphere}
$\CC\PP^1$ is homeomorphic to $S^2$, by an explicit map. Identify $\RR^3$ with $\CC \times \RR$ and let
\[
h \colon \CC\PP^1 \to S^2, \qquad h([z_0 : z_1]) = \frac{(2\bar{z}_0z_1, \ |z_1|^2 - |z_0|^2)}{|z_0|^2 + |z_1|^2}.
\]
Multiplying $(z_0, z_1)$ by $\lambda \in \CC^\times$ multiplies numerator and denominator by $|\lambda|^2$, so $h$ is well defined, and it is continuous since it is
induced by a continuous map on $\CC^2 \setminus \{0\}$. It takes values in $S^2$, because
$|2\bar{z}_0z_1|^2 + (|z_1|^2 - |z_0|^2)^2 = (|z_0|^2 + |z_1|^2)^2$. In the chart $U_0 = \{z_0 \neq 0\}$, with $w = z_1/z_0$,
\[
h([1 : w]) = \frac{(2w, \ |w|^2 - 1)}{|w|^2 + 1},
\]
which is the inverse of the stereographic projection from $N = (0, 1)$ (Example~\ref{general-topology:example-stereographic}, with $\RR^2 = \CC$). So $h$ maps
$U_0$ bijectively onto $S^2 \setminus \{N\}$, and the remaining point $[0 : 1]$ goes to $N$. Thus $h$ is a continuous bijection from the compact space $\CC\PP^1$ to
the Hausdorff space $S^2$, hence a homeomorphism: $\CC\PP^1$ is the one-point compactification of the chart $U_0 \cong \CC$, the \emph{Riemann sphere}. The
composite $S^3 \to \CC\PP^1 \to S^2$, $(z_0, z_1) \mapsto (2\bar{z}_0z_1, |z_1|^2 - |z_0|^2)$, is the Hopf map of Remark~\ref{homotopy:remark-higher-computations}.
For $\RR\PP^1$ the same construction gives $\RR\PP^1 \cong S^1$, the one-point compactification of $\RR$.
\end{example}

\section{Gluing}\label{topological-manifolds:section-gluing}

\begin{lemma}\label{topological-manifolds:lemma-gluing}
Let $X$ and $Y$ be spaces, $U \subseteq X$ and $V \subseteq Y$ open, and $f \colon U \to V$ a homeomorphism. Let $Z = X \cup_fY$ be the quotient of $X \sqcup Y$ by $u \sim f(u)$ for $u \in U$.
\begin{enumerate}
\item $X$ and $Y$ map homeomorphically onto open subsets of $Z$ covering $Z$.
\item If $X$ and $Y$ are Hausdorff, then $Z$ is Hausdorff if and only if the graph $\Gamma_f = \{(u, f(u))\}$ is closed in $X \times Y$.
\item If $X$ and $Y$ are $n$-manifolds and $\Gamma_f$ is closed, then $Z$ is an $n$-manifold.
\end{enumerate}
\end{lemma}

\begin{proof}
(1) The saturation of an open $A \subseteq X$ is $A \sqcup f(A \cap U)$, which is open, so the quotient map $q$ is open, and it is injective on $X$ and on $Y$. (2) Two points of $Z$ coming from $X$ are separated by images of
disjoint open subsets of $X$, and similarly for $Y$. The remaining case is $q(x) \neq q(y)$ with $x \in X \setminus U$, $y \in Y \setminus V$; then $(x, y) \notin \Gamma_f$. If $\Gamma_f$ is closed, there are open $A \ni x$, $B \ni y$ with
$(A \times B) \cap \Gamma_f = \emptyset$, that is $f(A \cap U) \cap B = \emptyset$, and $q(A)$, $q(B)$ are disjoint neighbourhoods. Conversely, if $(x, y) \in \bar\Gamma_f \setminus \Gamma_f$, then every pair of neighbourhoods $A \ni x$,
$B \ni y$ has $f(A \cap U) \cap B \neq \emptyset$, so $q(A) \cap q(B) \neq \emptyset$; and $q(x) \neq q(y)$ (if $x \in U$ then $y \neq f(x)$ since $(x, y) \notin \Gamma_f$; if $x \notin U$ and $y \notin V$ they are not identified; and $x \notin U$,
$y \in V$ is also not identified). (3) Locally Euclidean by (1), Hausdorff by (2), and second countable as the union of two open second countable subspaces.
\end{proof}

\begin{counterexample}\label{topological-manifolds:counterexample-gluing}
The line with two origins is $\RR \cup_f\RR$ with $U = V = \RR \setminus \{0\}$ and $f = \mathrm{id}$. The graph of $f$ is the diagonal minus the origin, whose closure in $\RR \times \RR$ contains $(0, 0)$; in accordance with
Lemma~\ref{topological-manifolds:lemma-gluing}(2), the result is not Hausdorff (Counterexample~\ref{topological-manifolds:counterexample-not-hausdorff}).
\end{counterexample}

\begin{definition}\label{topological-manifolds:definition-connected-sum}
Let $M$ and $N$ be connected $n$-manifolds, $n \geq 2$, and let $\varphi \colon U \to \RR^n$, $\psi \colon V \to \RR^n$ be charts onto $\RR^n$. Let $X = M \setminus \varphi^{-1}(\bar B_1)$ and $Y = N \setminus \psi^{-1}(\bar B_1)$, where
$\bar B_r = \{|x| \leq r\}$, and let $f$ be the homeomorphism from $\varphi^{-1}(\{1 < |x| < 2\}) \subseteq X$ onto $\psi^{-1}(\{1 < |y| < 2\}) \subseteq Y$ given in coordinates by $x \mapsto (3 - |x|)x/|x|$, which exchanges the two
boundary spheres of the annulus. The \emph{connected sum} is $M \# N = X \cup_fY$.
\end{definition}

\begin{proposition}\label{topological-manifolds:proposition-connected-sum}
$M \# N$ is a connected $n$-manifold, compact if $M$ and $N$ are.
\end{proposition}

\begin{proof}
$X$ and $Y$ are $n$-manifolds by Proposition~\ref{topological-manifolds:proposition-constructions}(1). We check that the graph of $f$ is closed in $X \times Y$. Let $(x_k, f(x_k)) \to (x, y) \in X \times Y$ with
$1 < |\varphi(x_k)| < 2$; sequences suffice, as $X \times Y$ is first countable. The points $\varphi(x_k)$ lie in the compact set $\{1 \leq |w| \leq 2\}$, so a subsequence converges to some $w$ there, and then the corresponding
subsequence of $(x_k)$ converges to $\varphi^{-1}(w)$; by uniqueness of limits in the Hausdorff space $M$, $x = \varphi^{-1}(w)$, and $|w| > 1$ because $x \in X$. In the same way $y = \psi^{-1}(w')$ with $|w'| > 1$ a limit point of
$\psi(f(x_k))$, whose norms are $3 - |\varphi(x_k)|$; so $|w'| = 3 - |w|$, which forces $|w| < 2$. Hence $x$ lies in the domain of $f$, and $y = f(x)$ by continuity of $f$. By Lemma~\ref{topological-manifolds:lemma-gluing},
$M \# N$ is an $n$-manifold.

$X$ is connected: the set $E = \varphi^{-1}(\{|w| > 1\})$ is connected, being homeomorphic to $S^{n-1} \times (1, \infty)$ with $S^{n-1}$ path-connected for $n \geq 2$ (as the image of $\RR^n \setminus \{0\}$ under $w \mapsto w/|w|$). If $X = A \sqcup A'$
with $A$, $A'$ open, we may assume $E \subseteq A$; then $M = (A \cup \varphi^{-1}(B_2)) \sqcup A'$ is a decomposition into disjoint open sets, since the points of $\varphi^{-1}(B_2)$ in $X$ lie in $E$; as $M$ is connected, $A' = \emptyset$.
Similarly $Y$ is connected, and the images of $X$ and $Y$ in $M \# N$ meet, so $M \# N$ is connected. If $M$ and $N$ are compact, the closed subsets $M \setminus \varphi^{-1}(B_{3/2}) \subseteq X$ and
$N \setminus \psi^{-1}(B_{3/2}) \subseteq Y$ are compact, and their images cover $M \# N$, because a point of $X$ with $1 < |\varphi| < 3/2$ is identified with a point of $Y$ with $3/2 < |\psi| < 2$.
\end{proof}

\begin{remark}\label{topological-manifolds:remark-connected-sum}
The homeomorphism type of $M \# N$ can depend on the choice of charts, through orientations: for orientable manifolds the two choices of relative orientation can give different results. For surfaces, and more
generally up to these orientation choices, the result is independent of the charts, but the proof in the topological category uses deep results; in the smooth category it follows from the disc theorem of Palais.
We shall only use connected sums of surfaces, where $T^2 \# T^2$ is the closed orientable surface of genus $2$ and $\RR\PP^2 \# \RR\PP^2$ is the Klein bottle.
\end{remark}

\begin{theorem}[Classification of surfaces]\label{topological-manifolds:theorem-classification-surfaces}
Every compact connected $2$-manifold is homeomorphic to exactly one of: the sphere $S^2$, a connected sum $T^2 \# \cdots \# T^2$ of $g \geq 1$ tori, or a connected sum $\RR\PP^2 \# \cdots \# \RR\PP^2$ of $k \geq 1$ projective planes.
\end{theorem}

This is \cite[Chapter 12]{munkres-topology}. The surfaces in the list are distinguished by their fundamental groups (their abelianisations are $\ZZ^{2g}$ and $\ZZ^{k-1} \oplus \ZZ/2$), and every compact surface is homeomorphic to one
obtained by identifying the edges of a polygon in pairs, which is the route to the proof. In dimension $1$ the classification is simpler: every connected $1$-manifold is homeomorphic to $\RR$ or to $S^1$
\cite[Appendix]{milnor-differential-viewpoint}.

\section{Exercises}\label{topological-manifolds:section-exercises}

\begin{exercise}\label{topological-manifolds:exercise-quotient-group}
Let $N$ be a normal subgroup of a topological group $G$. Show that $G/N$ with the quotient topology is a topological group, and that it is Hausdorff if and only if $N$ is closed.
\end{exercise}

\begin{solution}
The quotient map $\pi \colon G \to G/N$ is open, by the argument of Proposition~\ref{topological-manifolds:proposition-orbit-space}(1) for the action of $N$ by right translation. So $\pi \times \pi$ is open and surjective, hence a
quotient map, and multiplication $G/N \times G/N \to G/N$ is continuous because its composite with $\pi \times \pi$ is. Inversion is similar. By Proposition~\ref{topological-manifolds:proposition-topological-groups}(5), $G/N$ is
Hausdorff if and only if $\{eN\}$ is closed, that is, if and only if $\pi^{-1}(eN) = N$ is closed (the complement of $\{eN\}$ is open if and only if its preimage $G \setminus N$ is).
\end{solution}

\begin{exercise}\label{topological-manifolds:exercise-compact-second-countable}
Show that a compact Hausdorff space that is locally Euclidean is second countable, hence a manifold.
\end{exercise}

\begin{solution}
Finitely many chart domains $U_1, \ldots, U_k$ cover the space, and each is second countable, being homeomorphic to an open subset of $\RR^n$. The union of countable bases of the $U_i$ is a countable basis.
\end{solution}

\begin{exercise}\label{topological-manifolds:exercise-mobius}
Show that the Möbius band $B$ of Example~\ref{topological-manifolds:example-boundary} is the quotient of $\RR \times [-1, 1]$ by the free proper action of $\ZZ$ given by $n \cdot (x, t) = (x + n, (-1)^nt)$, and deduce that $B \setminus \partial B$ is a
$2$-manifold whose fundamental group is $\ZZ$.
\end{exercise}

\begin{solution}
Every orbit meets $[0, 1] \times [-1, 1]$, in one point or, for $x \in \ZZ$, in the two points $(0, t)$ and $(1, -t)$; the induced continuous bijection from the band to the orbit space is a homeomorphism since the band is compact and the
orbit space Hausdorff (Theorem~\ref{topological-manifolds:theorem-proper-free}, which applies because $K \cap (K + n) = \emptyset$ for compact $K$ and large $|n|$). Restricted to $\RR \times (-1, 1)$, which is simply connected and a
$2$-manifold, Theorem~\ref{topological-manifolds:theorem-covering-action} gives $\pi_1(B \setminus \partial B) \cong \ZZ$.
\end{solution}

\begin{exercise}\label{topological-manifolds:exercise-hopf}
Let $\ZZ$ act on $\RR^n \setminus \{0\}$ by $k \cdot x = 2^kx$. Show that the action is free and proper, and that the quotient is homeomorphic to $S^{n-1} \times S^1$.
\end{exercise}

\begin{solution}
The map $\RR^n \setminus \{0\} \to S^{n-1} \times \RR$, $x \mapsto (x/|x|, \log_2|x|)$, is a homeomorphism carrying the action to $k \cdot (u, s) = (u, s + k)$, which is free and proper on $S^{n-1} \times \RR$ with quotient
$S^{n-1} \times (\RR/\ZZ) \cong S^{n-1} \times S^1$ (Example~\ref{topological-manifolds:example-orbit-spaces}(1)). The action on $S^{n-1} \times \RR$ is proper
because a compact set lies in some $S^{n-1} \times [-r, r]$, which meets its translate by $k$ only if $|k| \leq 2r$. For $n = 2$ the quotient is a torus. For $n \geq 3$, $\RR^n \setminus \{0\} \cong S^{n-1} \times \RR$ is simply
connected, $S^{n-1}$ being simply connected (Chapter~\ref{homotopy}), so the quotient has $\pi_1 \cong \ZZ$ by Theorem~\ref{topological-manifolds:theorem-covering-action}.
\end{solution}
